\chapter{Details}\label{ch:Details}

\section{Werkzeuge}
\begin{mytheorem}[geometrische Summe]\label{mytheorem:geometrischeSumme}
	Es sei $q \neq 0$ eine reelle Zahl und $n$ eine natürliche Zahl. Dann gilt
	\begin{align*}
		\sum_{k = 0}^n q^k = q^0 + q^1 + \ldots + q^n = \frac{q^{n+1}-1}{q-1}.
	\end{align*}
	\begin{myproof}
		Wir beweisen die Aussage durch vollständige Induktion.
		\begin{itemize}
			\item Für $n = 0$ gilt die Aussage, da
			      \begin{align*}
				      \sum_{k = 0}^0 q^k = q^0 = 1 = \frac{q^{0+1}-1}{q-1}.
			      \end{align*}
			\item Die Aussage gelte nun für eine natürliche Zahl $n$. Wir zeigen, dass sie auch für $n+1$ gilt.
			      \begin{align*}
				      \sum_{k = 0}^{n+1} q^k = q^{n+1} + \sum_{k = 0}^n q^k = q^{n+1} + \frac{q^{n+1}-1}{q-1} = \frac{q^{n+2}-1}{q-1}.
			      \end{align*}
		\end{itemize}
	\end{myproof}
\end{mytheorem}


\begin{mytheorem}[verallgemeinerte geometrische Summe]\label{mytheorem:verallgemeinertegeometrischeSumme}
	Es sei $q \neq 0$ eine reelle Zahl und $n$ und $m\leq n$ natürliche Zahlen. Dann gilt
	\begin{align*}
		\sum_{k = m}^n q^k = q^m + q^{m+1} + \ldots + q^n = \frac{q^{n+1}-q^m}{q-1}.
	\end{align*}
	\begin{myproof}
		Unter Verwendung von \cref{mytheorem:geometrischeSumme} finden wir
		\begin{align*}
			\sum_{k = m}^n q^k =
			\sum_{k = 0}^n q^k - \sum_{k = 0}^{m-1} q^k =
			\frac{q^{n+1}-1}{q-1} - \frac{q^m - 1}{q-1} =
			\frac{q^{n+1}-q^m}{q-1}.
		\end{align*}
	\end{myproof}
\end{mytheorem}


